\begin{figure}[H]
\centering
\begin{tikzpicture}[
    scale=0.85, transform shape,
    % --- STYLES ---
    action/.style={
        rectangle, rounded corners=0.3cm,
        draw=black, fill=white,
        minimum width=3.4cm, 
        minimum height=1cm,
        align=center, font=\small,
        inner sep=4pt
    },
    decision/.style={
        diamond, draw=black, fill=white,
        aspect=2, inner sep=2pt,
        font=\footnotesize, minimum width=2.5cm,
        align=center
    },
    start/.style={circle, fill=black, minimum size=0.5cm},
    end/.style={circle, draw=black, double, minimum size=0.5cm, inner sep=1pt},
    control/.style={->, thick, >=Stealth, rounded corners=5pt},
    swimlane/.style={thick, black!60}
]

% ===== 1. COORDINATES & HEADERS =====
\def\xManager{0} 
\def\xAcc{8} 

\node[font=\bfseries\large] at (\xManager, 1) {Quản lý};
\node[font=\bfseries\large] at (\xAcc, 1) {Kế toán};

% Đường kẻ phân cách 2 cột
\draw[swimlane] (4, 1.5) -- (4, -16.5);

% ===== 2. MANAGER LANE NODES (Phần Đầu) =====
\node[start] (start) at (\xManager, 0) {};

\node[action, below=0.6cm of start] (m1) {Thông báo chi tiết khấu trừ\\ \& số tiền hoàn/thu thêm};

\node[action, below=0.8cm of m1] (m2) {Xác nhận đồng ý trả phòng\\ theo kết quả đối soát};

\node[decision, below=0.8cm of m2] (dec1) {Khách cần\\ thanh toán thêm?};

% Nhánh Có (Thu thêm)
\node[action, below=1cm of dec1] (m3_yes) {Thông báo thu thêm\\ cho Kế toán};


% ===== 3. ACCOUNTANT LANE NODES (Luồng Thu Thêm) =====
\node[action] (a1_yes) at (\xAcc, 0 |- m3_yes) {Hướng dẫn khách\\ thanh toán};

\node[action, below=0.8cm of a1_yes] (a2) {Xác nhận thanh toán\\ của khách};

\node[decision, below=0.8cm of a2] (dec2) {Thanh toán\\ thành công?};

% Nhánh thay thế (Thất bại) rẽ phải để tạo vòng lặp
\node[action, right=1.2cm of dec2] (a_fail) {Yêu cầu khách\\ thanh toán lại};


% ===== 4. KẾT HỢP HAI LUỒNG (Bên dưới) =====
% Nhánh Không (Hoàn cọc) - Kế toán
\node[action, below=1.5cm of dec2] (a1_no) {Thống nhất phương thức\\ thanh toán với khách};

% Nhánh Không (Hoàn cọc) - Quản lý
\node[action] (m3_no) at (\xManager, 0 |- a1_no) {Thông báo hoàn cọc\\ cho Kế toán};

% Hội tụ cả 2 luồng về UC4-3
\node[action, below=1cm of a1_no] (a3) {Thực hiện UC4-3};

\node[end, below=0.8cm of a3] (end_main) {};
\fill (end_main.center) circle (0.15cm);


% ===== 5. CONNECTIONS =====

% -- Luồng Quản lý Ban Đầu --
\draw[control] (start) -- (m1);
\draw[control] (m1) -- (m2);
\draw[control] (m2) -- (dec1);

% -- Rẽ nhánh 1 (Thu thêm hay Hoàn cọc) --
\draw[control] (dec1) -- node[right, font=\footnotesize] {Có} (m3_yes);

% Nhánh Không đi vòng qua mép trái xuống khối m3_no
\draw[control] (dec1.west) -- node[above, font=\footnotesize] {Không} ++(-2.2,0) |- (m3_no.west);

% -- Chuyển giao sang Kế toán --
\draw[control] (m3_yes) -- (a1_yes);
\draw[control] (m3_no) -- (a1_no);

% -- Luồng Thu Thêm của Kế toán --
\draw[control] (a1_yes) -- (a2);
\draw[control] (a2) -- (dec2);

% Vòng lặp yêu cầu thanh toán lại
\draw[control] (dec2) -- node[above, font=\footnotesize] {Không} (a_fail);
\draw[control] (a_fail.north) |- ([yshift=5pt]a2.east);

% -- Mũi tên hội tụ về UC4-3 --
% Luồng Hoàn cọc đâm thẳng xuống
\draw[control] (a1_no) -- (a3);

% Luồng Thu Thêm (Có) đi vòng qua bên phải để không cắt ngang khối a1_no
\draw[control] (dec2.south) -- node[left, font=\footnotesize] {Có} ++(0,-0.6) -| ([xshift=2.5cm]a3.east) -- (a3.east);

\draw[control] (a3) -- (end_main);

\end{tikzpicture}
\caption{Activity Diagram UC4-2: Đối soát và thanh toán khi trả phòng}
\end{figure}
